\documentclass[12pt]{article}
\usepackage{amsmath,amssymb,amsfonts,amsthm}
\usepackage[margin=1in]{geometry}

\begin{document}

\begin{center}
{\Large\bfseries Chapter 4, Problem 5}\\[6pt]
Byron \& Fuller, \textit{Mathematics of Classical and Quantum Physics}
\end{center}

\bigskip

\textbf{Problem.} Prove that the eigenvalues of any anti-self-adjoint matrix are imaginary.

\bigskip
\textbf{Solution.}

Let $A$ be anti-self-adjoint, so $A^\dagger = -A$. Suppose $Ax = \lambda x$ for some nonzero vector $x$.

Take the inner product of both sides with $x$:
\[
(x, Ax) = \lambda (x, x) = \lambda \|x\|^2.
\]

Now consider $(x, Ax)$. Using the property of the adjoint:
\[
(x, Ax) = (A^\dagger x, x) = (-Ax, x) = -\overline{(x, Ax)}.
\]

So $(x, Ax) = -\overline{(x, Ax)}$, which means $(x, Ax)$ is purely imaginary (a complex number that equals the negative of its conjugate is purely imaginary).

Since $\|x\|^2 > 0$ (as $x \neq 0$) and $(x, Ax) = \lambda \|x\|^2$, we have:
\[
\lambda = \frac{(x, Ax)}{\|x\|^2},
\]
which is purely imaginary.

Therefore all eigenvalues of an anti-self-adjoint matrix are purely imaginary. \qed

\end{document}
